\section{Introduction}
\label{sec:intro}

Reconciliation-based delivery gave cloud operations a powerful invariant:
the running system converges to a declared desired state, and divergence---%
\emph{configuration drift}---is detected and repaired
continuously~\cite{opengitops,beetz2022,argocd,fluxcd,burns2016}. The invariant
is so useful that it has quietly become a proxy for a much stronger claim. When
the reconciler reports \textsc{Synced} and the drift detector reports clean,
operators, auditors, and dashboards read: \emph{the system is in the state we
approved}. This paper is about the gap between those two statements.

Consider a payments service whose deployment was approved on March~3rd: the
manifest matched revision \texttt{rev-42}, satisfied policy version
$\pi_7$, carried a service-owner approval bound to image digest
\texttt{sha256:aaa\ldots}, and traced to ticket \texttt{TKT-99}. Months later
the cluster still matches its Git manifest byte for byte. Meanwhile: the
security team superseded $\pi_7$ with $\pi_8$, which the running configuration
violates; an emergency exception granted during an incident expired
after four hours and has now been silently in force for six weeks; the registry
re-tagged \texttt{svc:1.4.2} so the pods run a digest no approval ever
referenced; and a well-meaning rollback reverted Git to \texttt{rev-37},
content whose approval belongs to a different change. Configuration drift:
zero. One or more legitimacy conditions that covered the deployment on
March~3rd no longer hold. We call this condition \emph{governance drift}: the
operational state has diverged not from its manifest but from the state,
policy context, authorization basis, and intent \emph{under which it was
approved}.

The distinction is not terminological. The desired/observed
configuration-drift signal implemented by reconcilers and commonplace IaC
drift tools is a \emph{two-place, memoryless} relation between two
configurations and is therefore decidable by comparison
\cite{awsconfig,driftctl}. Configuration-management programs can retain
approved historical baselines (\cref{sec:related}), but their decision object
remains configuration-scoped. Governance drift is a
\emph{three-place, history-anchored} relation among present state, evolving
context, and an immutable admitted basis; we establish the separation in two complementary
forms: constructed, tool-realizable counterexamples, and a
definitional-dependency result showing that no detector whose inputs are
configuration pairs---however sophisticated its diffing---can decide
governance consistency, because the deciding facts are not functions of those
inputs (\cref{prop:impossible}). Worse, reconciliation can materialize an
unauthorized desired-state change and then erase the transient
configuration-drift signal, leaving intent drift: the rollback scenario above
is converged into content outside the historically authorized change lineage
(\cref{sec:taxonomy}).

Our executed study makes the separation concrete. We implement a
governance-drift detector evaluated at cumulative evidence tiers---%
configuration comparison; plus policy re-evaluation; plus authorization
records with validity semantics; plus artifact lineage; plus environment
inventory against recorded assumptions---against twelve injected drift
scenarios (including compound drift and evidence decay) and a no-drift control containing both
runtime churn and \emph{benign governance churn} (satisfied policy
revisions, approved updates, hygienic exceptions), under a counterfactual
scoring rule that credits information rather than alarm volume. The result is
a strict detectability staircase (\cref{tab:matrix}): normalized
configuration comparison detects exactly one scenario of twelve; each evidence
tier detects precisely the scenarios whose deciding facts it carries; the
governance tiers achieve zero false alarms \emph{earned against} governance
churn. A deliberately naive unnormalized comparator is retained only as a
scoring diagnostic, not as a production baseline.

\subsection{Research Questions}
\begin{description}[leftmargin=3.2em,itemsep=1pt,topsep=2pt]
\item[RQ0:] Can divergence between approved intent and operational reality be
formally represented and measured as governance drift?
\item[RQ1:] Which forms of governance drift can arise and be distinguished in
cloud-native systems?
\item[RQ2:] Can governance drift be detected using existing operational,
governance, and control-plane evidence streams?
\item[RQ3:] How quickly does governance drift become visible?
\item[RQ4:] Which forms of drift are detectable through configuration
comparison alone, and which require authorization or evidence context?
\end{description}

RQ0 is answered by the formal model and the combined analytic and empirical
results. RQ1 is answered by the taxonomy and scenario catalog
(\cref{sec:taxonomy,sec:study}); RQ2 and RQ4 by the dependency result (\cref{prop:impossible}), the tiered architecture (\cref{sec:architecture}),
the executed detectability study (\cref{sec:study}), and the bounded live
laboratory (\cref{sec:lab}); RQ3 by the study's event-time measurements and
the laboratory's separated actuation, onset-to-detection, end-to-end, and
time-to-evidence observations across three polling cadences.

\subsection{Contributions and Epistemic Status}
\begin{enumerate}[label=\textbf{C\arabic*.},leftmargin=2.2em,itemsep=1.5pt,topsep=2pt]
\item \textbf{Taxonomy} (conceptual): formal definitions separating
five substantive components---configuration, policy, authorization, intent,
and environment---plus transversal epistemic evidence drift, with governance
drift as their umbrella---and with the explicit
commitment that governance drift must not, and provably does not, reduce to
configuration drift (\cref{sec:taxonomy}).
\item \textbf{Formal model} (conceptual/analytic): an immutable sequence of
atomic approval snapshots plus separate activation and supersession records,
stable unit identity, admissible cross-stream temporal cuts, observed state $\Gobs(t)$, time-indexed governance contexts
$\Pol(t), \Auth(t), \Int(t)$, and a component-wise governance-divergence vector
$\GDsub(u,\theta)$ plus an observability mask $\ObsMask_{\mu}(u,\theta;T)$, with a reasoned rejection of a primitive scalar distance
(\cref{sec:model}).
\item \textbf{Separation results} (analytic): constructed counterexamples and
a definitional-dependency proposition establishing that configuration
equality does not imply governance consistency and that no configuration-pair
detector decides it---deliberately elementary, and billed as such
(\cref{sec:taxonomy,sec:model}).
\item \textbf{Detection architecture} (design): a vendor-neutral tiered
architecture mapping each component to a sufficient evidence tier, together
with indistinguishable-world constructions proving tier minimality
(\cref{sec:architecture,prop:tier-minimality}).
\item \textbf{Executed detector study} (empirical, closed-world): twelve
scenarios $\times$ six detector tiers $\times$ twenty seeds, with
counterfactual scoring, full-vector ground truth, three compound scenarios,
exact-set and Hamming scoring, a separately implemented composite comparator without the
admitted-basis join, and a control containing
runtime \emph{and} benign governance churn; all code and data packaged,
matrix generated programmatically (\cref{sec:study}). Its role is precisely
scoped: it validates that the reference semantics realize the model's
dependency structure and quantifies noise behavior; it is not field evidence,
and it makes no claim about any real cluster.
\item \textbf{Repeated Kubernetes/GitOps laboratory} (empirical, executed): a
version-pinned Kind/Flux/Kyverno stack, local Git and OCI services, approved-state
recorder, tiered evaluator, and executable injections for S1--S9. In one
hundred eighty controlled injections (20 per scenario), three randomized-phase
polling cadences produced 538/540 detections and 100\% conditional exact-set
accuracy for provisional class-set classifications under the laboratory's
sequential-snapshot assumption. Twenty steady-state benign windows produced no alarms in 543 polls;
a separate transition-inclusive campaign began three isolated observers before
each benign mutation and retained all 777 polls, including seven expected
configuration-convergence polls and 36 fail-safe epistemic warning polls but
no policy, authorization, intent, or environment drift; all 60 trajectories
ended with a provisionally consistent classification
(\cref{sec:lab}). A secondary S10--S12 extension reached
the exact provisional class set in 45/45 cadence observations; six separate 300-s benign windows
added 4,680 steady-state polls without substantive or epistemic alarms. The execution
demonstrates bounded realizability and behavior within the reported laboratory conditions, not natural occurrence,
prevalence, production effectiveness, stream-loss reliability, or transfer.
\item \textbf{Boundary campaigns} (empirical, executed): a nine-episode audit
with 27 separate observer processes and 261 raw polls; nine controlled profiles
through a two-hop localhost TCP evidence path; an in-memory batched-join
benchmark through 1,000 units; and a live Kubernetes object-path campaign that
completed 10- and 50-Deployment targets while enforcing declared stop rules.
The process-isolated audit reached the exact target in 27/27 correlated
trajectories. Cause-start-to-first-exact completion had median 4.465~s
(nearest-rank P95 9.764), and cause-start-to-two-poll exact persistence had
median 9.522~s (P95 19.811). The latter is only a second consecutive exact poll,
not watermark-qualified Stable-VCL. A bounded Argo CD/Gatekeeper replication
(\cref{sec:lab}) produced 15/15 exact projected singleton classifications over
declared evaluated components. S1 configuration and S3 policy used native
second-stack paths; S4 reused the shared artifact adapter. Intent and
environment were not evaluated, so the result is neither a complete formal
vector evaluation nor a cross-stack equivalence test.
The separate temporal-cut tests validate the admissibility contract rather
than retroactively making the sequential live snapshots globally coherent.
These campaigns
test traceability, fail-safe transport semantics, decision-core scaling, live
fan-out, and limited component-path transfer separately; none is presented as
an equivalence, production-capacity, or reliability estimate.
\end{enumerate}

Established results are cited; analytic claims are proved in the model; the
closed-world and live-stack empirical claims are separated explicitly; and
the repeated live-stack protocol is reported with its negative controls and
misses rather than presented as production evidence.

\subsection{Scope and Non-Goals}
We address drift \emph{after} an approved deployment; admission-time
enforcement~\cite{k8sadmission,opa,kyverno,binauthz} is complementary and
assumed imperfect (break-glass paths exist by design). We propose no product;
the detector logic we release is a reference semantics, not a system. And we
take no position on \emph{which} governance controls organizations should
impose---the model makes drift relative to whatever was approved under
whatever policy; the choice of policy is the organization's.
